\section*{第 \ref{ch:b2:mantle-melting} 章——地幔熔融}

\begin{solution}{exo:b2:mantle-melting:1}
(a) 橄榄石超过 \qty{90}{\percent}：\omterm{def:b2:mantle-melting:peridotites}{纯橄岩}。(b) 橄榄石与斜方辉石，
不含单斜辉石：\omterm{def:b2:mantle-melting:peridotites}{方辉橄榄岩}。(c) 橄榄石超过 \qty{40}{\percent}，并含两种辉石：
\omterm{def:b2:mantle-melting:peridotites}{二辉橄榄岩}。(d) 橄榄石少于 \qty{40}{\percent}：辉石岩
（橄榄网辉石岩，同时含两种辉石）。(e) $57/98 = \qty{58.2}{\percent}$，
$28/98 = \qty{28.6}{\percent}$，$13/98 = \qty{13.3}{\percent}$：\omterm{def:b2:mantle-melting:peridotites}{二辉橄榄岩}。
\end{solution}

\begin{solution}{exo:b2:mantle-melting:2}
(a) 在 \qty{60}{\percent} 钙长石处液相线为 \qty{1390}{\celsius}，位于钙长石相区：
钙长石首先结晶。(b) 透辉石出现在
低共熔点 \qty{1270}{\celsius}，此时液相含钙长石 \qty{42}{\percent}。
(c) 两种矿物任一混合物的第一批液相都是低共熔液，
含钙长石 \qty{42}{\percent}，温度为 \qty{1270}{\celsius}。
\end{solution}

\begin{solution}{exo:b2:mantle-melting:3}
(a) $\num{4.0e-5}\times1623/(3300\times1000) = \qty{1.97e-8}{K/Pa} =
\qty{19.7}{K/GPa}$。(b) $\num{e9}/(3300\times9.81) = \qty{3.09e4}{m}$，或
每吉帕 \qty{30.9}{km}；$19.7/30.9 = \qty{0.64}{K/km}$，约为
近地表 \qty{30}{K/km} 的五十分之一，那里热量以传导方式传递，
而非由物质输运。
\end{solution}

\begin{solution}{exo:b2:mantle-melting:4}
(a) $T_{\mathrm s}(0) = \qty{1085.7}{\celsius}$；$T_{\mathrm s}(1) = 1085.7 +
132.9 - 5.1 = \qty{1213.5}{\celsius}$; $T_{\mathrm s}(2) = 1085.7 + 265.8 - 20.4 =
\qty{1331.1}{\celsius}$; $T_{\mathrm s}(3) = 1085.7 + 398.7 - 45.9 =
\qty{1438.5}{\celsius}$。(b) $132.9 - 2\times5.1\times2 = \qty{112.5}{K/GPa}$。
(c) 在 \qty{3}{GPa} 时，\qty{1400}{\celsius} 低于固相线，
\qty{1438.5}{\celsius}；沿绝热线升至 \qty{2}{GPa} 时，温度仅降低约
\qty{20}{\celsius}，降至约 \qty{1380}{\celsius}，而固相线降至
\qty{1331.1}{\celsius}：岩石已处于固相线之上，部分发生熔融。
\end{solution}

\begin{solution}{exo:b2:mantle-melting:5}
(a) $1350 + (1277.5 - 1350)\times(30 - 20)/(40 - 20) = \qty{1314}{\celsius}$。
(b) 该成分位于低共熔点的透辉石一侧：透辉石。
(c) 液相以 \qty{42}{\percent} 钙长石到达低共熔点；钙长石
仍留在液相中，因此液相分数为 $30/42 = 0.714$，于是 $1 - 0.714 =
\qty{29}{\percent}$ 的体系以透辉石形式结晶。(d) 全部
结晶：\qty{70}{\percent} 透辉石与 \qty{30}{\percent} 钙长石，即
整个体系的成分。
\end{solution}

\begin{solution}{exo:b2:mantle-melting:6}
(a) 镁橄榄石：$60.08/(2\times40.30 + 60.08) = 60.08/140.68 =
\qty{42.71}{\percent}$ \ce{SiO2}；顽火辉石：$60.08/100.38 = \qty{59.85}{\percent}$。
(b) $0.975\times59.85 + 2.5 = \qty{60.85}{\percent}$。(c) 顽火辉石位于镁橄榄石与熔体之间，
镁橄榄石的分数为 $(60.85 - 59.85)/(60.85 -
42.71) = 1.00/18.14 = \qty{5.5}{\percent}$，即鲍文与安德森给出的值。(d) 熔体
中的二氧化硅含量高于其来源的斜方辉石，而残余物
中橄榄石增多。
\end{solution}

\begin{solution}{exo:b2:mantle-melting:7}
(a) $5.1P^2 - 112.9P + 214.3 = 0$：$P_0 = (112.9 - \sqrt{12746.4 -
4371.7})/10.2 = (112.9 - 91.51)/10.2 = \qty{2.10}{GPa}$，即 $2.097\times30.9 =
\qty{64.8}{km}$。(b) $5.1P^2 - 112.9P + 364.3 = 0$：$P_0 = (112.9 -
72.90)/10.2 = \qty{3.92}{GPa}$，即 \qty{121}{km}。(c) $T_{\mathrm s}'(2.10) = 132.9
- 21.4 = \qty{111.5}{K/GPa}$，因此 $\Delta P_0 = 10/(111.5 - 20) = \qty{0.109}{GPa}$
每 \qty{10}{\celsius}，或 \qty{3.4}{km}。
\end{solution}

\begin{solution}{exo:b2:mantle-melting:8}
(a) $T_{\mathrm s} = \qty{1331.1}{\celsius}$，$T_{\mathrm l} = 1475.0 + 160.0 -
12.8 = \qty{1622.2}{\celsius}$; $T = 1331.1 + 0.10\times291.1 =
\qty{1360.2}{\celsius}$，或 \qty{1633}{K}。(b) $(\partial T/\partial P)_F = 112.5
+ 0.10\times(67.2 - 112.5) = \qty{108.0}{K/GPa}$; $(\partial T/\partial F)_P =
T_{\mathrm l} - T_{\mathrm s} = \qty{291.1}{K}$。(c) $c_p/T = 0.612$；取 $\alpha/\rho
= \qty{12.1}{J/(kg.K)}$ 每吉帕，分子为 $0.612\times108.0 - 12.1
= 54.0$，分母为 $300 + 0.612\times291.1 = 478.2$：$-\mathrm dF/\mathrm
dP = 0.113$，即每吉帕 \qty{11.3}{\percent}。(d) $54.0/178.2 = 0.303$，
每吉帕 \qty{30}{\percent}：潜热使产率
降低 2.7 倍。
\end{solution}

\begin{solution}{exo:b2:mantle-melting:9}
(a) $43\times1 = \qty{43}{\celsius}$；$43\times5^{0.75} = 43\times3.34 =
\qty{144}{\celsius}$; $43\times13^{0.75} = 43\times6.85 = \qty{294}{\celsius}$.
(b) $1213.5 - 294 = \qty{919}{\celsius}$。(c) \qty{125}{ppm} 即
\qty{0.0125}{\percent}，因此 $X = 0.0125/0.01 = \qty{1.25}{\percent}$，于是 $\Delta T =
43\times1.25^{0.75} = 43\times1.18 = \qty{51}{\celsius}$.
\end{solution}

\begin{solution}{exo:b2:mantle-melting:10}
(a) 底部宽 $2H$，地幔以 $U$ 的速度穿过它：$2UH$，即每单位洋脊长度、
每单位时间。每一侧跨越高度 $H$，并以水平速度穿过
$U$：每侧 $UH$，共 $2UH$。若侧面倾角为 $\theta$，则底宽为
$2H/\tan\theta$，流入量为 $2UH/\tan\theta$，两者相等仅当
$\tan\theta = 1$：本模型中倾角由质量守恒确定。(b) 产率恒定时
\omterm{def:b2:mantle-melting:fraction}{熔体分数}在顶点处为 $F_{\max} = \Pi P_0 = \Pi\rho
g z_0$，且 $H = z_0$：$h_c = F_{\max}z_0/2 = \Pi\rho g z_0^2/2$。(c) 每年
新增 \qty{50}{mm} 板块，$\qty{0.05}{m}\times\qty{1000}{m}\times\qty{7100}{m} =
\qty{3.6e5}{m^3}$ 的地壳，每千米洋脊；平均\omterm{def:b2:mantle-melting:fraction}{熔体分数} \qty{12}{\percent} 时
它来自 $\num{3.6e5}/0.12 = \qty{3.0e6}{m^3}$ 的地幔。
\end{solution}

\begin{solution}{exo:b2:mantle-melting:11}
(a) 从单位质量液体中析出质量为 $t$ 的 A 后，液体含 A 为 $0.6 -
t$，B 为 0.2，C 为 0.2，质量为 $1 - t$：$\big((0.6 - t), 0.2,
0.2\big)/(1 - t)$，B/C 保持为 1。(b) 沿\omterm{def:b2:mantle-melting:ternary}{共结线}，$(0.40 - 0.15s, 0.60 -
0.30s, 0.45s)$；令 B = C 得 $0.60 - 0.30s = 0.45s$，$s = 0.8$：$L_1$ 中 A 为 28，B
为 36，C 为 36。于是 $0.2/(1 - t) = 0.36$，$t = 0.44$：\qty{44}{\percent} 的体系已
以 A 的形式结晶。(c) C 全部在液相中：$f = 20/45 = 0.44$，
\qty{44}{\percent} 仍为液相。(d) A：$0.60 - 0.444\times0.25 = 0.49$；B：$0.20 -
0.444\times0.30 = 0.07$；加上液相后，$49 + 7 + 44 = \qty{100}{\percent}$。
\end{solution}

\begin{solution}{exo:b2:mantle-melting:12}
(a) $X_{\mathrm{sat}}(1) = 12 + 1 = \qty{13}{\percent}$，$\Delta T =
\qty{294}{\celsius}$，饱和固相线 $1213.5 - 294 = \qty{919}{\celsius}$；
$X_{\mathrm{sat}}(2) = 12\times1.516 + 2 = \qty{20.2}{\percent}$, $\Delta T =
43\times20.2^{0.75} = 43\times9.53 = \qty{410}{\celsius}$，饱和固相线 $1331.1
- 410 = \qty{921}{\celsius}$。(b) 熔体在较高压力下溶解更多的水，
较大的降幅抵消了干固相线的 \qty{118}{\celsius} 升高。
(c) 固相线必须下降 $1331.1 - 1100 = \qty{231}{\celsius}$：$X =
(231/43)^{1/0.75} = 5.37^{1.333} = \qty{9.4}{\percent}$ 的水含量，低于
饱和时的 \qty{20.2}{\percent}，因此地幔楔可以熔融。
\end{solution}

\begin{solution}{pb:b2:mantle-melting:1}
\textbf{1.} $1350 + 273 = \qty{1623}{K}$。
\textbf{2.} $\num{4.0e-5}\times1623/(3300\times1000) = \qty{1.97e-8}{K/Pa} =
\qty{19.7}{K/GPa}$，四舍五入为 \qty{20}{K/GPa}。
\textbf{3.} $\num{e9}/(3300\times9.81) = \qty{30.9}{km}$; $20/30.9 =
\qty{0.65}{K/km}$。
\textbf{4.} $1350 + 20P = 1085.7 + 132.9P - 5.1P^2$，即 $5.1P^2 - 112.9P +
264.3 = 0$：$P_0 = (112.9 - \sqrt{12746.4 - 5391.7})/10.2 = (112.9 - 85.76)/10.2 =
\qty{2.66}{GPa}$。
\textbf{5.} $2.66\times30.9 = \qty{82}{km}$；$T = 1350 + 20\times2.66 =
\qty{1403}{\celsius}$。
\textbf{6.} $1085.7 + 132.9\times2.66 - 5.1\times2.66^2 = 1085.7 + 353.5 - 36.1 =
\qty{1403}{\celsius}$。
\textbf{7.} $T_{\mathrm s} = 1085.7 + 199.4 - 11.5 = \qty{1273.6}{\celsius}$；
$T_{\mathrm l} = 1475.0 + 120.0 - 7.2 = \qty{1587.8}{\celsius}$；差值为
\qty{314.2}{K}。
\textbf{8.} $1273.6 + 0.10\times314.2 = \qty{1305.0}{\celsius}$，\qty{1578}{K}。
\textbf{9.} $T_{\mathrm s}' = 132.9 - 15.3 = \qty{117.6}{K/GPa}$，$T_{\mathrm l}' =
80.0 - 9.6 = \qty{70.4}{K/GPa}$：$(\partial T/\partial P)_F = 117.6 +
0.10\times(70.4 - 117.6) = \qty{112.9}{K/GPa}$；$(\partial T/\partial F)_P =
\qty{314.2}{K}$。
\textbf{10.} $c_p/T = 1000/1578 = 0.634$；分子 $0.634\times112.9 - 12.1 = 71.6
- 12.1 = 59.5$；分母 $300 + 0.634\times314.2 = 499.2$：$0.119$，即
\qty{11.9}{\percent}，每吉帕。
\textbf{11.} 取拟合函数 $F = T'^{1.5}$，其中 $T'$ 为从固相线到液相线的相对位置，温度必须随 $F$ 急剧升高，才能产生
最初百分之一的熔体，因此 $(\partial T/\partial F)_P$ 在固相线附近较大，
熔融开始得很慢。$F_{\max} = 0.09\times2.66 = 0.24$，\qty{24}{\percent}。
\textbf{12.} $0.13/0.50 = 0.26$：单斜辉石在
\qty{26}{\percent} 熔融时耗尽。在 \qty{24}{\percent} 时，轴线下方的残余体仍保留少量
单斜辉石，为贫单斜辉石的\omterm{def:b2:mantle-melting:peridotites}{二辉橄榄岩}，位于\omterm{def:b2:mantle-melting:peridotites}{方辉橄榄岩}的边缘。
\textbf{13.} $h_c = 0.24\times82/2 = \qty{9.8}{km}$。
\textbf{14.} $F_{\max}/2 = \qty{12}{\percent}$。
\textbf{15.} \qty{9.8}{km} 比 \qty{7.1}{km} 多 \qty{38}{\percent}，超出
正常地壳的范围。
\textbf{16.} 产率恒定时，最深处的地幔产生的熔体过多；
拟合模型使最初的熔体缓慢生成，$F(z)$ 在起始处保持较低，并且
$\int F\,\mathrm dz$ 较小。
\textbf{17.} $F_{\max} = \Pi P_0$ 和 $H = 30.9P_0$（km）：$h_c = \Pi P_0\times30.9P_0/2
= kP_0^2$，其中 $k = 0.09\times30.9/2 = \qty{1.39}{km/GPa^2}$。
\textbf{18.} $P_0 = \sqrt{7.1/1.39} = \sqrt{5.11} = \qty{2.26}{GPa}$。
\textbf{19.} $2.26\times30.9 = \qty{70}{km}$。
\textbf{20.} $T_{\mathrm s}(2.26) = 1085.7 + 300.4 - 26.0 = \qty{1360.0}{\celsius}$，
则 $T_p = 1360.0 - 20\times2.26 = 1360.0 - 45.2 = \qty{1315}{\celsius}$。
\textbf{21.} 对于 \qty{6.3}{km}：$P_0 = \sqrt{4.53} = \qty{2.13}{GPa}$，$T_{\mathrm s}
= \qty{1345.6}{\celsius}$，$T_p = \qty{1303}{\celsius}$。对于 \qty{7.9}{km}：$P_0 =
\sqrt{5.68} = \qty{2.38}{GPa}$，$T_{\mathrm s} = \qty{1373.1}{\celsius}$，$T_p =
\qty{1326}{\celsius}$。地壳的离散允许\omterm{def:b2:mantle-melting:tp}{地幔位温}在 $1315 \pm
\qty{12}{\celsius}$ 范围内波动。
\textbf{22.} $P_0 = \sqrt{20/1.39} = \sqrt{14.39} = \qty{3.79}{GPa}$；$T_{\mathrm s} =
1085.7 + 503.7 - 73.3 = \qty{1516.1}{\celsius}$；$T_p = 1516.1 - 75.8 =
\qty{1440}{\celsius}$，高出 $1440 - 1315 = \qty{125}{\celsius}$。
\textbf{23.} 两者依据的证据不同。熔岩记录的是最热的地幔，
位于地幔柱中心；地壳平均了所有穿过熔融区的地幔熔体，
其中大部分温度较低。简单模型还保持
产率恒定，尽管单斜辉石在 \qty{26}{\percent} 处耗尽，
而此处 $\Pi P_0 = 0.09\times3.79 = 0.34$：真实的热地幔熔融较少，需要
更高的温度才能生成相同的地壳。
\textbf{24.} $T_{\mathrm s}'(2.26) = 132.9 - 23.1 = \qty{109.8}{K/GPa}$；$\Delta P_0 =
10/(109.8 - 20) = \qty{0.111}{GPa}$；由于 $h_c \propto P_0^2$，$\Delta h_c/h_c =
2\times0.111/2.26 = \qty{9.8}{\percent}$，约每 \qty{10}{\celsius} 增加十分之一。
\textbf{25.} 在极慢速的洋脊下，地幔上升得极为缓慢，以至于通过
传导向冷的表面散热；它偏离绝热线，其熔融
在更深处停止，三角形顶部（$F$ 最大处）被切去，地壳
变薄。
\textbf{26.} 正常洋脊的地壳给出的\omterm{def:b2:mantle-melting:tp}{地幔位温}为
\textbf{\qty[text-series-to-math=true]{1315}{\celsius}}，误差约在
\qty{12}{\celsius} 以内：比矿物物理学给出的 \qty{1350}{\celsius} 低 \qty{35}{\celsius}，
该值的不确定度与之大致相当，并且很好地落在 1280 至
\qty{1400}{\celsius} 之间（由熔岩给出）。三种相互独立的温度计都指示地幔
约为 \qty{1300}{\celsius} 至 \qty{1350}{\celsius}。
\end{solution}
